\chapter[Sermon 69. The Three Former Angels Pour Out Their Vials upon the Antichristians, and All the Ungodly.]{The Three Former Angels Pour Out Their Vials upon the Antichristians, and All the Ungodly.}
\label{sermon:069}
\markright{Sermon 69. The Three Former Angels Pour Out Their Vials upon the ...}

And I heard a great voice out of the temple saying to the angels: Go your ways, pour out your vials of wrath upon the Earth. And the first went, and poured out his vial on the earth, and there fell a noisome sore upon the men which had the mark of the beast, and upon them that worshipped his image. And the second angel poured out his vial on the sea, and it turned as it were into the blood of a dead man: and every living thing died in the sea. And the third angel poured out his vial upon the rivers and fountains of waters, and they turned to blood, and I heard an angel of the waters saying: Lord which art and wast, thou art righteous and holy, because thou hast given such judgements: for they shed the blood of saints, and prophets, and therefore thou hast given them the blood to drink: for they are worthy. And I heard another angel out of the altar saying: Even so, Lord God almighty, true and righteous are thy judgements.

After speaking generally about God’s righteous judgments, he now proceeds particularly through the number seven and declares at length the plagues of God, which he also inflicts in this world upon the wicked, but especially upon those who oppose Christ. This passage corresponds to the same, or at least has many similarities to, the passages in Moses’s book of Exodus from chapter 7 through chapter 12. For in those chapters are described the ten plagues of God, with which, for sin, he afflicted King Pharaoh and the entire realm of Egypt. These plagues are summarized in fine verses by Dr. Musculus, our worthy Godfather.

These plagues are also explained in Psalm 150. In chapter 15 of Exodus, the Lord says: “If you diligently hear the voice of your God, and do what is right in his sight, and keep all his statutes, I will not send upon you any disease that I sent upon the Egyptians, for I am the Lord, healing you.” We learn, therefore, from the account of God’s plagues, to fear God and to walk in his commandments. It is not contrary to this saying of God that we read how Job and other holy men, walking in the commandments of God, were afflicted with grievous diseases. For these are private [illegible] are not chiefly inflicted for sin, but for the exercise of faith and the increase of virtues.

Men for the most part ascribe the causes of plagues to the stars and to other matters, and therefore do not turn to the Lord, who is striking them in response to a most wicked life. But we are taught by the treatise of Moses, which we cited from Exodus, and by this present discussion of John, that God himself punishes the sins and wickedness of men, although he uses the service of men and elements. To these, as to the proximate causes, men attribute the evils they receive, which they justly suffer from God for their sins. For this cause, a voice is now heard, not from the air or from the earth, but from the Temple of the Lord, true, just, and holy, commanding the angels to pour out their vials upon the heads of men. The wicked are therefore plagued by God himself. But a vial is nothing other than the just judgment of God, or vengeance deserved by men. Angels pour out their vials whenever men are punished with plagues through means appointed by God. And that voice which is heard from the temple is great, for no man can resist God, nor infringe his decree. When he commands, all creatures obey.

\index{Antichrist}But while this first angel, executor of God’s judgment, unleashes his plague upon men, a terrible sore botch fell upon them. This plague corresponds to the sixth plague of Egypt. And that botch signifies a canker, a fistula, and swelling sores or boils, but chiefly the pox of Jude, which others call the disease of Naples, some the French pox, and some the Spanish; truly, for that in the war of Naples (which was waged by the French and Spaniards in the year of our Lord 1494) it first appeared in the camp of prostitutes, infecting the army. Mainardus the physician discusses this at length. But however diverse and venomous sores do infect many grievously, yet the French pox chiefly corrupts the abbeys of monks and nuns, and colleges of priests, above others. For they giving themselves to most filthy fornication, do abhor and detest in others holy matrimony, and therefore receive thereof the reward of their iniquity. Therefore it is said here expressly that the Antichrist should be vexed with this disease, or rather punished. You shall find some whose face is eaten with this disease. All whoremongers and adulterers for the most part are troubled with this plague. Job also, the excellent servant of God, was covered with sores and boils, but by the singular counsel of God, as I touched upon before. Therefore it is no marvel, though sometimes very good men, free from the uncleanness of whoredom, are also infected with this disease.

The second angel poured out his vial on the sea, and immediately the sea became like the blood of a dead man, corrupt and lifeless. All that lived in the sea died. The sea is always in motion and changeable; therefore, it rightly signifies the world, or unconstant people in the world, who, because of their sins, are afflicted with the plague and die in great numbers. In the words there is the figure of synecdoche, where every living soul is said to die. This second plague corresponds to the fifth plague of Egypt. Under this plague we understand all kinds of pestilences and plagues. Hezekiah also suffered from the plague, as do many godly men, but by God’s unique counsel.

The third angel poured his vial upon the rivers and fountains of waters, which were immediately turned into blood. This corresponds to the first plague of Egypt. The Egyptians had drowned in the Nile the newborn infants and had oppressed the innocent Israelites; therefore, they were worthy to drink of the Nile, for it became blood.

\index{Augustine, Saint}\index{Zechariah (Book of)}\index{Ezekiel (Book of)}Water otherwise in the Scripture signifies doctrine, as in Ezekiel and Zechariah. Therefore, the rivers and fountains of water signify ecclesiastical preachers and rulers, whom God has given to the people for defense and relief. Certainly, Saint Peter calls false prophets dry wells. In 2 Peter 2:17, we shall hear that by waters are understood people. This, then, is God’s plague: rulers of the people and preachers of peace have become the instigators and leaders of rebellion and war, in which they fall and kill one another, shedding the blood of the saints. And although in wars the godly are also afflicted, the Lord knows how to requite their suffering and ease their sorrows. Saint Augustine settles this matter at length in the first book of Christian doctrine. But if we look upon the tumultuous history of Italy, France, Germany, and Hungary, and other realms that glory in being called Christian, we shall find them to have been blazing firebrands of war, which by right ought to have been the rulers of peace.

\index{Turks}\index{Papacy}\index{Rome}As the Lord says in the Gospel, a prophet must not die anywhere except in Jerusalem, so no war must be instigated except by the Popes of Rome, bishops, and prelates. I will only recount a few examples. Pope Gregory II, through sedition, expelled Emperor Leo Isauricus from Italy. Pope Stephen brought King Philip of France into Italy against the Lombards. The same did Charlemagne at the urging of Pope Leo III, driving him clean out of Italy, having slain many of them with the sword. Pope Gregory VII, a most wicked man, stirred up Peter, King of Hungary, to war with Emperor Henry IV, entangled all of Italy and Germany with wars, and drove Henry to fight many battles, which were not light. Urban II, of that name, stirred up war in both the East and West and all other parts of the world he called holy, attempting to recover Jerusalem. This war was long, cruel, great, and bloody, such as cannot be found in all the world. What Alexander III accomplished against Frederick Barbarossa, and how he raised up all of Italy against him, the histories tell. And while Frederick II warred in the Holy Land, Gregory IX took Naples from him. Here, the Abbot of Ursperge observes that such wickedness should be committed by a Pope. Great factions arose in Italy, of the Guelphs and Ghibellines, through the means and instigations of the Popes. Clement IV brought in the French army, under the leadership of King Charles, into the kingdom of Naples, and deposed Conrad, Duke of Swabia, from his inheritance, and caused Frederick, Duke of Austria, and Conrad to be slain together with many thousands of Germans. Pope John XXII armed Frederick, Duke of Austria, and Leopold against Emperor Louis IV of the house of Bavaria. Boniface VIII commanded King Albert, Duke of Austria, to bear hostile banners against King Philip of France. As Martin V stirred up a grievous war against the Bohemians. Eugenius IV betrayed Ladislaus, King of Poland and Hungary, to Amurath, the great Turk, to be vanquished and slain through treason, sending his legate Julian Caesarini, a cardinal, about the practice, who also perished in that disastrous defeat. This recalls the saying in Virgil: “Even the soothsayer himself is slain.” Pope Sixtus IV sent to the powerful Swiss nation a red scarf with a bull attached, granting large indulgences to those who would fight for the church of Rome. Julius II, through much bloodshed of the Swiss, began to expel the French king from Italy, which at length Leo X accomplished, receiving Emperor Charles V, whose son still rules in Italy. Clement VII began to oppress him again, but death thwarted his endeavors. Paul III joined the forces of Italy with Charles V, and warred against the Germans for obedience denied to the See of Rome and the Gospel received. In this war, Philip, the Landgrave of Hesse, and John Frederick, Duke of Saxony, an Elector Prince, were taken. Great villainy and cruelty were wrought by the soldiers in Germany. Pope Julius III began to parley with the Frenchmen, and stirring up the war of Parma and Mirandula, brought the Frenchmen to Siena. There arose a most grievous war by sea and land, both in France and Italy, and also in Germany, which continues to this day: the princes and people tear one another asunder, they drink their blood abundantly, while nevertheless, in the meantime, they persecute Christ’s church most grievously. The Lord send peace.

\index{John, Saint}\index{Daniel (Book of)}And now where the godly might marvel why God so suffers the world to be shaken and troubled with mutual wars, the Angels prevent the marveling and complaint, and show not only the cause, but also praise the justice of God in these judgments. And he brings in two Angels, as meet and sufficient witnesses of this business. The one he makes ruler of waters, the other speaking out of the altar. He seems here to follow Daniel, who also in the tenth chapter, says that Angels as governors were set to rule over provinces. Not that God does not work and govern all things in waters and in all elements and regions; but for that he uses the travel of Angels as his ministers. But where the Papists gather hereof that Saints rule over elements, diseases, limbs, cities, and every part in man, it is foolish and superstitious, and smelling of idolatry. For the manner of Angels and of blessed souls is clean diverse, moreover, the Scripture attributeth unto them far other things than it does to these. You shall read nothing of the blessed souls as having anything to do with men here on Earth in the whole Scripture. But in sundry places of the Scriptures, you shall read that Angels are set to be men's keepers, and to serve them with divers ministries. Again, you read not that the godly have for this cause given any godly honor to the Angels; no, we shall hear in this book how Saint John would have worshipped an Angel, but was prohibited of the Angel once or twice. Moreover, here the Angel renders a reason why the water is turned into blood, and commends here in God’s justice. For turning his talk unto God, “Thou art indeed, O Lord which art, and which wast,” and so forth. He pronounces him righteous, as he that will do no man any wrong, and therefore calls him also holy. In the meantime he signifies his everlastingness, and that he gives being unto all things, where he says, “which art, and which wast,” and so forth. Of this phrase of speech is spoken in the first chapter. And the true righteousness gives to every one his. Therefore the Angel says, “Therefore Lord thou art righteous, and declarest thy righteousness to the world, in that thou hast given them blood to drink, which have shed the blood of the Prophets,” that is, of preachers, for preaching of the truth. And not their blood only, but have shed also the blood of thy holy faithful—I mean, whom for the true professing of the faith they have vexed, and at last slain. Therefore are they worthy that they themselves should again drink the blood of them and theirs; that is, should fall by mutual wars, tumults, and slaughters, verily before recited.

These things are confirmed by another angel who speaks from the altar, and not without cause from the altar. For we heard before in chapter 6 that under the altar the souls of those who have been killed cry out and say, “How long will you not avenge our blood?” Therefore, now a voice is uttered from the altar so that we should understand that God does not forget the blood of his saints, but avenges it in just and due season. Now, here, by the way, the omnipotence of God is commended, so that the ungodly may understand that in the time of affliction and vengeance, there will be no power able to resist the Almighty. To him alone be glory. Amen.
